\section{Problem Setup}

\paragraph{Complex Tasks.} In this work our goal is to build a generalist agentic system capable of solving \emph{complex tasks} across a variety of domains. We define a task as {\em complex} if it requires, or significantly benefits from, a process involving planning, acting, observing, and reflecting, potentially multiple times. Acting refers to more than generating tokens, such as executing code, using tools, or interacting in an environment. Observing, in this context, provides information that was previously unavailable or unknowable. A task is defined by an input, a desired output and an evaluation function to compare the desired output to any candidate output. The input consists of a well-specified textual description and an optional arbitrary set of file attachments which may include images, dataset files, audio clips among other things. For example, the input task description could be \textit{``fact-check each claim in the attached PDF as correct or incorrect"} with a PDF file as an attachment. The desired output consists either of a textual answer (possibly representing a structured object), or a specific state of the environment to reach. In the fact-checking example, the output might be a string labeling each fact as correct or not, e.g., \textit{``claim 1: correct, claim 2: incorrect, ..."}.
Here, the evaluation function might simply determine whether the desired output and the proposed answer match exactly.


\paragraph{Agentic Systems.} 
To complete a task, assume a {\em computer} which can be partially observed and operated to complete the task. The computer constitutes the \emph{environment}.
An agentic system can take as input the task description, and any related attachments that are present on the computer environment. The system is allowed to do arbitrary processing to complete the task, but must complete it within a time budget (e.g., 25 mins).
For instance, on the computer, the autonomous system can execute Python code, navigate the web using a browser, download files locally, among other actions from its \emph{action space}. The system's ability to take action in, and potentially modify, both the local and web environments is why we refer to the system as \emph{agentic}. After completing the task, the system returns a text answer, and a trace of its observations and steps along the way. The final state of the environment is also captured in sufficient detail to run the task evaluation. 
Note that this setting can be described as a Partially Observable Markov Decision Process, similar to formalizations used by prior work~\cite{sodhi2024stepstackedllmpolicies}.
Nex, we describe \team, our multi-agent system that can autonomously solve complex tasks. 